\section{Leaderboard Competition in Practice}
We now investigate whether capability profiles are differentiated in practice and if performance converges over time. 
We first study capability profile differentiation by looking at whether developers exhibit distinct strengths (RQ1) and whether developers improve more in categories  where they lag competitors (RQ2). 
We then study convergence in performance by examining whether the competitive frontier becomes more crowded over time (RQ3). 
Together, these analyses ask whether the patterns generated by our theoretical model appear in observed model competition rather than directly testing the assumptions of our theoretical model. 

\subsection{Dataset Details}
We rely on two sources of data: leaderboards from LM Arena and usage data from OpenRouter.
\paragraph{Leaderboards}
We use LM Arena to construct leaderboards and better understand how capability profiles evolve over time. 
We construct a panel covering 422 models from 59 organizations and 28 categories, using the ``lmarena-ai/leaderboard-dataset'' from March 2024 through September 2026. 
Each snapshot contains 28 category-level Elo ratings, along with rating variances and vote counts. 
Of these 28 categories, we focus on 8 (coding, math, hard prompts, instruction following, creative writing, longer query, multi-turn, and English) because these categories are consistently available since 2024. 

\paragraph{OpenRouter}
To measure realized model usage, we use daily OpenRouter traffic data from January 2025 through September 2026. The data cover 408 models over 620 days.
We aggregate token shares to the model and developer levels to study how changes in the evaluated model frontier correspond to changes in realized model usage.
We detail how we match data from OpenRouter up with data from LM Arena in Appendix~\ref{app:matching_details}. 

\subsection{RQ1: Do developers retain distinct capability profiles?}
\begin{figure*}
    \includegraphics[width=0.657\textwidth]{figures/residual_cosine_kde.pdf}
    \includegraphics[width=0.343\textwidth]{figures/residual_cosine_heatmap.pdf}
    \caption{Developers exhibit distinct capability profiles overall, with clusters of similar relative strengths. At each snapshot, we represent each developer by its highest-rated model and compare residualized capability profiles. Left: distribution of pairwise cosine similarities across developers, compared with a permutation null that randomly reassigns capabilities across developers. Right: pairwise similarities reveal clusters of developers with similar relative strengths.}
    \label{fig:residual_cosine_kde}
\end{figure*}

We first examine whether developers' capability profiles are differentiated, motivated by Section~\ref{sec:monoculture}. 
Our goal is to isolate which categories each developer performs relatively well on, accounting for both category difficulty  (e.g., math being easier than English) and dominant developers (e.g., Anthropic being broadly better). 
However, directly applying the metrics from Section~\ref{sec:monoculture} to Elo scores is inappropriate because Elo scores are identified only up to an additive shift. 
Instead, at each timestep, we select the highest rated model for each developer and compare capability profiles. 
To account for differences in overall developer strength and category difficulty, we fit $\theta_{i,k} = \delta + \eta_i + \gamma_k + r_{i,k}$ for each time step separately. 
Here, differences in $r_{j,k}$ capture differences in capability profiles, while $\eta_i$ captures differences in overall performance. 
This metric captures similarities in the capability profiles for different developers, rather than directly measuring the theoretical cosine similarity from Section~\ref{sec:monoculture}.  

In Figure~\ref{fig:residual_cosine_kde} (left), we plot the distribution of $\cos(r_{i,\cdot},r_{i',\cdot})$ between pairs of developers.
We find little evidence of a common capability profile across all developers, as the observed similarities are close to a null model (see Appendix~\ref{app:statistical_rq1} for details and statistical tests). 
However, in Figure~\ref{fig:residual_cosine_kde} (right), we do find that clusters of developers tend to be similar. 
For example, models from Google and Meta are stronger in creative writing and multi-turn interactions and weaker in coding and math. 
These clusters cut across developer types; for example, models from OpenAI and Anthropic are in a different cluster than those from Google and Meta. 

\begin{figure*}[t]
    \includegraphics[width=0.5\textwidth]{figures/residual_timeseries.pdf}
    \includegraphics[width=0.5\textwidth]{figures/residual_norm.pdf}
    \caption{Some developer-specific capability advantages persist over time, even as the overall magnitude of specialization changes. Left: most developer-category residuals are noisy, but a few persist across time. For example, Anthropic's score in English and its score in instruction following hold in sign across time steps. Right: magnitude of residuals over time.}
    \label{fig:residual_distribution}
\end{figure*}

We next ask whether these relative strengths persist over time. 
Figure~\ref{fig:residual_distribution} (left) shows that many developer--category residuals fluctuate around zero, but some exhibit stable advantages or disadvantages. 
In particular, we find that 29 of the 104 developer--category pairs have the same sign in every quarter. 
For example, Anthropic persistently underperforms in English, while Cohere overperforms in creative writing. 
Figure~\ref{fig:residual_distribution} (right) plots the  L2 norm of each developer's residual vector, \(\lVert r_{i,\cdot,t} \rVert_2\), which captures the magnitude of category-specific specialization. 
We see that this magnitude first increases then decreases over time, indicating that developers' capability profiles become less specialized after a certain time. 
This pattern holds both for the 14 developers that appear in every quarter and the full set of developers, including those entering in 2025 and 2026. 
Together, these results suggest that developers retain some persistent differences in their relative strengths, even as the overall magnitude of specialization decreases. 
While the magnitude of specialization decreases in later periods, we find little evidence of a uniform capability monoculture.

\subsection{RQ2: Do developers improve more in capabilities where they lag competitors?}
We investigate how differences in capability profiles relate to subsequent improvement by looking at how developers improve or catch up to competitors. 
This examines an empirical analogue of the catch-up condition in Proposition~\ref{prop:AsymptoticBunching}.
Specifically, we relate a developer's improvement during new model releases, \(\Delta\theta_{i,k,t}\), to its pre-release gap relative to the best rival \(g_{i,k,t} = \theta_{\mathrm{leader},k,t} - \theta_{i,k,t}\).
We fit $\Delta\theta_{i,k,t} = \eta_{i,k} + \gamma_{k,t} + \beta g_{i,k,t} + \varepsilon_{i,k,t}$, where $t$ is the calendar quarter of the release, and $\eta$ and $\gamma$ represent fixed effects (details in Appendix~\ref{app:statistical_rq2}). 
We estimate $\hat{\beta}$, which captures whether larger competitive gaps predict subsequent improvement, for the 12 developers with sufficiently many model releases (more details in Appendix~\ref{app:statistical_rq2}).

\begin{figure*}[t]
    \includegraphics[width=0.5\textwidth]{figures/gap_vs_dtheta.pdf}
    \includegraphics[width=0.5\textwidth]{figures/beta_forest.pdf}
    \caption{Developers catch up more in categories where they lag further behind the leader. Left: relationship between competitive gaps and subsequent improvements, where $\hat{\beta}$ denotes the slope. Right: developer-specific estimates of the catch-up coefficient, with 95\% confidence intervals.}
    \label{fig:evalchasing_scatter}
\end{figure*}

Figure~\ref{fig:evalchasing_scatter} shows that developers catch up more in categories where they lagged the frontier, with an estimated coefficient of $\hat{\beta}=0.48$. 
Intuitively, a 10-point larger competitive gap is associated with 4.8 points more subsequent improvement, conditional on the fixed effects.
This relationship holds across most developers, as $\hat{\beta}>0$ for 10 out of 12 developers, with the 95\% confidence interval excluding zero. 
Here, the standard error is similar when clustering by either release event ($0.117$) or developer ($0.113$).

\begin{figure*}[t]
    \centering 
    \includegraphics[width=\textwidth]{figures/beta_by_level.pdf}
    \caption{The catch-up relationship remains positive after controlling for overall model strength. Left: estimates of \(\hat{\beta}\) after adding a linear control for overall model strength, \(\theta_{i,\cdot}\). Right: estimates of \(\hat{\beta}\) across quartiles of overall model strength.}
    \label{fig:score_control}
\end{figure*}

To better understand the factors underlying this pattern, we investigate whether this improvement reflects differences in overall model strength, developers' category-specific weaknesses, or the strength of competitors. 
Figure~\ref{fig:score_control} accounts for overall baseline model performance, \(\theta_{i,\cdot}\), in two ways: by directly controlling for \(\theta_{i,\cdot}\) in the regression (left), and by estimating \(\hat{\beta}\) separately across quartiles of \(\theta_{i,\cdot}\) (right).
In both cases, \(\hat{\beta}\) remains positive, though the catch-up relationship is largest among models with lower baseline performance.
We additionally decompose the frontier gap into a developer's own category-specific performance and the performance of its strongest rival in Appendix~\ref{app:strong_competitors}.
Our results suggest that developers particularly catch up in categories where they themselves are weak, rather than simply in categories with unusually strong rivals.
Together, these results suggest that the observed improvement pattern is primarily associated with developers' category-specific weaknesses, rather than their overall performance or unusually strong competitors.

We conduct additional robustness checks in Appendix~\ref{app:confounds} by accounting for release year, regression to the mean, and comparator choice. 
The positive relationship persists across alternative specifications, while our regression-to-the-mean analysis suggests that measurement noise alone does not explain the observed association.
Together, RQ1 and RQ2 reveal that developers retain distinct relative strengths across capabilities, even as they tend to improve more in areas where they lag their competitors. 
We next examine how performance gaps between models have changed over time. 

\subsection{RQ3: Are models converging in performance?}
\begin{figure*}[t]
    \includegraphics[width=\textwidth]{figures/top_rank_gaps.pdf}
    \caption{Performance gaps narrow across the competitive frontier. We measure gaps between the leader and the 10th-, 20th-, and 50th-ranked models, averaged across categories.}
    \label{fig:near_top}
\end{figure*}

\begin{figure*}[t]
  \includegraphics[width=0.33\textwidth]{figures/top_distribution.pdf}
  \includegraphics[width=0.33\textwidth]{figures/usage_within_points.pdf}
  \includegraphics[width=0.33\textwidth]{figures/effective_number.pdf}
    \caption{Models become closer in performance over time, alongside increasingly dispersed usage. Left: models in 2026 lie closer to the performance frontier than in earlier years. Middle: OpenRouter token usage is concentrated among models near the Arena leader. Right: the effective number of models and developers used on OpenRouter has increased over time.}
    \label{fig:top_distro}
\end{figure*}

Having found that developers tend to improve more in capabilities where they lag competitors (RQ2), we next examine whether performance gaps between competing models have narrowed over time.
We do so by measuring the performance gap between the leader and lower-ranked models, which is the empirical analog of $\mathbf{g}^{t}$ from Proposition~\ref{prop:AsymptoticBunching}, averaged across categories. 
Figure~\ref{fig:near_top} shows increasing crowding; between July 2024 and September 2026, the absolute gaps between the leader and the 10th-, 20th-, and 50th-ranked models decrease by 53\%, 61\%, and 66\%, respectively.  
This pattern is not driven solely by developers releasing multiple models: the number of distinct developers within 100 and 200 points of the leader increases by 2.4$\times$ and 1.4$\times$, respectively. 
We further investigate this in Figure~\ref{fig:top_distro} (left), where we show that models in 2026 lie closer to the frontier than models from 2025. 
These results demonstrate performance convergence for an increasingly large set of models over time. 

We next investigate whether this crowded frontier is accompanied by changes in model usage. 
We do so by matching data between Arena and OpenRouter, covering 71.3\% of overall OpenRouter token volume, and 76.7\% average monthly token volume (see Appendix~\ref{app:matching_details}).
Figure~\ref{fig:top_distro} shows that OpenRouter usage is concentrated among near-frontier models; models within 150 points of the Arena leader account for 70\% of observed tokens. 
At the same time, OpenRouter usage is becoming increasingly dispersed across models and developers. 
Between January 2025 and August 2026, the effective number of models and developers--defined as the inverse Herfindahl index of token shares--used on OpenRouter increases by 2.4\(\times\) and 2.7\(\times\), respectively.
Increasing competition near the performance frontier coincides with a broader distribution of realized model usage.
Together, our empirical results reveal convergence in performance without uniform monoculture in capability profiles. 
Developers retain persistent differences in their relative strengths (RQ1), even as improvements are concentrated in capabilities where they lag competitors (RQ2) and the performance frontier becomes increasingly crowded (RQ3). 